Developing the Hermite tokenization branch taught me to define what a model's action representation can express. Separately, I adapted OpenPI for a CR5 loop with WebSocket policy serving and real-robot execution. These efforts gave me experience with representation and physical policy integration, but they did not constitute an end-to-end agent spanning reasoning, model orchestration, cloud infrastructure, and embodied action. The Hong Kong Polytechnic University (PolyU)'s MSc in Agentic AI Systems targets precisely this systems-level gap.

Subject to departmental offering in my intake, COMP5583 Agentic Computing, Systems, and Applications would give me a framework for coordinating an agent's decisions and tools; COMP5585 Cloud Computing and Hyperscale AI Infrastructure would expose the deployment layer beyond my laboratory WebSocket service; and ME5301 Embodied Robot Intelligence would connect those abstractions back to sensing and control. If I selected the optional project route and received approval for the six-credit COMP5934 Agentic AI Systems Project, I could test one bounded agent pipeline from decision trace to physical command without overstating my prior CR5 work as a large-model or agent deployment. My habit of separating codec, policy, service, and hardware evidence would support careful team evaluation and prepare me to engineer accountable embodied agents for robotics R\&D.
